
Masinnägemise süsteemi arenduseks leiti sobivad komponendid, mis lõpptulemusena toimisid omavahel heas koostöös, tagades süsteemi eduka ning nõuetekohase toimimise riistvara ja tarkvara näol.

Riistvaraliselt valitud juhtkontrolleriks Nvidia Jetson AGX Xavier, mis suudab edukalt mitme sisendiga, suuremaid ning täpseid närvivõrgu mudeleid jooksutada. Samuti võimaldas ka antud kontrolleri riistvara suure kasuteguriga mudelite kiirendamist, mis aitas oluliselt masinnägemise rakenduse tööl kaameralt saadud kaadrite töötlemise kiirusele kaasa. Samuti aitas ka täpsematele objektide tuvastamisele kaasa kaamerateks valitud Arducam Fisheye moodul võimaldas edastada juhtkontrollerile nõuetekohase resolutsiooniga informatsiooni.

Tarkvaraliselt oli olulisel kohal kogu rakenduse konteinerdamine, mis võimaldas rakenduse arendamist ning toodangkeskkonnas käitamist hõlpsamaks teha, kuna sellisel meetodil oli rakendust võimalik käitada erinevatel arvutisüsteemidel ilma muudatusteta süsteemi enda tarkvaras.

Masinnägemise tarkvara olulisim komponent oli objektide tuvastuseks valitud meetod ehk masinnägemise jaoks loodud närvivõrgumudelid. Nende seast valiti YOLOv5 Ultralyticsi poolt, mis oli parim võrreldes teiste saadavalolevate mudelitega oma ühendatud objektide tuvastuse ning nende sildistamise arhitektuuri ja parima kiiruse ning täpsuse suhte poolest. Samuti oli see üks vähestest saadavalolevatest mudelitest, mida oli võimalik edukalt riistvaraliselt kiirendada.

Kõige olulisemad tulemused masinnägemise mudelite katsetes

\begin{itemize}
\item
  Masinnägemise mudelite täpsus kasvab YOLOv5m / YOLOv5l mudelini
\item
  Keskmine täpsuse muutus optimeeritud ja optimeerimata mudelite võrdluses oli väike \({1,7 \cdot 10}^{- 3}\)
\item
  Keskmine mudelite kiiruse kasv pärast optimeerimist oli 3,5 kordne
\item
  Parima optimeeritud mudeli kiiruse kasv võrreldes optimeerimata versiooniga oli 4,3 kordne
\end{itemize}

Masinnägemise rakenduse optimeerimisel saavutati sujuv ja kiirendatud rakenduse töö, kusjuures optimeerimise tulemusena vähenes juhtkontrolleri energiakasutus olulisel määral (26,8 \%) ning seetõttu olid ka töötemperatuurid madalamad. Samuti toimus selle tõttu ka andmete töötlus ja närvivõrgu töö oluliselt kiiremini ning kogu juhtkontrolleri riistvaralised kiirendid olid otstarbekalt kasutuses.
